\documentclass[10pt]{article}
\usepackage[top=0.77in,bottom=0.78in,left=0.88in,right=0.88in]{geometry}
\usepackage[cmex10]{amsmath}
\usepackage{amsthm}
\usepackage[hidelinks]{hyperref}
\DeclareFontShape{OMX}{cmex}{m}{n}{<->cmex10}{}
\DeclareMathSizes{10}{10}{7}{7}
\DeclareMathSizes{9}{10}{7}{7}
\linespread{0.96}
\newtheorem{theorem}{Theorem}
\newtheorem{lemma}[theorem]{Lemma}
\title{The strong Aharoni--Berger conjecture through five colours}
\author{Achita Anantachina}
\date{26 November 2026}
\begin{document}
\maketitle

\begin{abstract}
The strong Aharoni--Berger conjecture asks for a full rainbow matching among
\(n\) matchings of size \(n+1\) in a bipartite graph. We verify the conjecture
for \(n\leq5\) using three independently checked clausal refutations. One
refutation proves a mixed-size extension theorem for four colours. Applying
that theorem after deleting the endpoints of an edge leaves two endpoint
patterns for five colours. We specify the finite formulas and the reductions
that make their unsatisfiability sufficient.
\end{abstract}

\section{Statements}
Let \(M_0,\ldots,M_{n-1}\) be matchings between disjoint left and right vertex
sets. A full rainbow matching chooses exactly one edge from each \(M_c\), and
all selected left endpoints and all selected right endpoints must be
distinct. An edge can belong to several \(M_c\). The strong conjecture was
posed by Aharoni and Berger \cite{ab} and appears as Conjecture~2.7 in the
properly coloured bipartite multigraph formulation in \cite{dp}. Parallel
occurrences of the same edge in different colours encode exactly the same
instances.

\begin{theorem}\label{thm:five}
For \(1\leq n\leq5\), any \(n\) bipartite matchings with at least \(n+1\) edges
each have a full rainbow matching.
\end{theorem}

\begin{lemma}\label{lem:mixed}
Four bipartite matchings of sizes at least \(4,4,5,5\), respectively, have a
full rainbow matching.
\end{lemma}

Truncating each colour to its required size preserves any obstruction, so
we work with equality in these size thresholds. We first give the finite
formula for Lemma~\ref{lem:mixed}. For each side, introduce a Boolean
variable for equality of the endpoints of each unordered pair of edge
occurrences in different colours. Set equalities within each colour to
false. Add the three transitivity implications on every triple, simplifying
those that involve a fixed false literal. This describes two equivalence
relations with no repeated endpoint inside a colour. Conversely, using
the classes as vertices makes every satisfying assignment a valid
collection of matchings.

For every selection of one occurrence of each colour, require equality of
some pair of selected endpoints on at least one side. This adds one clause
per selection. To reduce symmetry, choose and fix a rainbow of the first
three colours, which exists greedily because their sizes are \(4,4,5\).
Only three vertices on either side are occupied, so at least two edges of
the last matching have a free left endpoint and at least two have a free
right endpoint. In an obstruction no edge has both endpoints free. Reorder
the last matching to put two edges of each type in four prescribed slots.
The formula enforces these free and occupied endpoints by units and
disjunctions. Its checked refutation proves Lemma~\ref{lem:mixed}.

\section{Reduction to endpoint patterns}
The case \(n=4\) follows from the lemma. The cases \(n=1,2\) follow by
greedy selection. For \(n=3\) we use the following elementary observation.

\begin{proof}[Auxiliary three-colour argument]
Write the third size-four matching as the diagonal on four indices.
If no rainbow exists, any edge \(i\to j\) of the first size-three matching
can be paired with an edge of the second, since it blocks at most two of
its three edges. Such a pair must hit all four diagonal edges. Therefore
every edge of the first two matchings is a loopless arc on these indices.
For each first-colour arc the second colour must supply an arc on the
complementary pair. Three arcs with distinct tails and heads on four
indices form a directed 3-cycle, a 2-cycle and a disjoint arc, or a directed
path with three arcs. A 3-cycle requires three second-colour arcs incident
with the fourth index, which is impossible. For a 2-cycle and a disjoint
arc, the latter requires a second-colour arc on the 2-cycle, whose reverse
in the first colour gives a rainbow with an unused diagonal. For the path
\(1\to2\to3\to4\), the required unordered pairs in the second colour are
\(34,14,12\). Reversing either end arc again gives an unused diagonal.
Otherwise the pair \(14\) must be oriented \(4\to1\), and together with
\(1\to2\) it misses diagonal index 3. This proves the auxiliary claim.
\end{proof}

Suppose five matchings, each of size six, form an obstruction. For each
edge \(e\) in colour \(c\), delete its endpoints. Every other colour loses at
most two edges. If at least two of those colours lost at most one edge,
Lemma~\ref{lem:mixed} would give a rainbow after deletion, and we could add
\(e\). Thus \(e\) must destroy two distinct edges in at least three of the
other four colours.

Every used endpoint therefore belongs to at least four colours. There are
30 incidences on each side, so each side has six or seven vertices. On a
side of size six every colour uses all vertices. On a side of size seven
there is one vertex missing each colour, and those five vertices are
distinct. The remaining two vertices are common to all five colours.
The edge condition also implies that an edge occurs in at most two
colours, or just one if either endpoint misses a colour.

If both sides have seven vertices, join only equal missing labels when
both endpoints miss a colour. Two different missing labels would cause
the edge to destroy at most one edge in two colours. Each colour uses four
labelled vertices per side and at most two can go to common vertices on
the opposite side. It therefore uses at least two equal-label edges. Each
can appear only once across colours, and there are only five such edges.
This contradicts the required ten uses.

\section{The remaining formulas}
By exchanging the sides if needed, there are six left vertices and either
six or seven right vertices. Label the right vertices arbitrarily in the
six-vertex case. In the seven-vertex case, label the missing vertices by
their colours \(0,\dots,4\), and the common vertices by 5 and 6. Relabel the
left vertices so that colour 0 is the fixed matching pairing them in order
with its eligible right vertices. The following clauses specify each
remaining formula.

\begin{enumerate}
\setlength{\itemsep}{0pt}
\item Six variables are fixed true for the prescribed colour-0 edges.
The other variables register the edges in the four remaining colours.
Each such colour has exactly one edge through each eligible row and column.
\item Each cell has multiplicity at most two, including occurrences in
colour 0. If its right vertex misses a colour, the cap is one.
\item For each fixed colour-0 edge and each choice of four more edges with
distinct rows and columns, require at least one of the four variables to
be false.
\end{enumerate}

An actual obstruction gives a satisfying assignment in one of these two
formulas. Both have checked refutations, which completes the proof of
Theorem~\ref{thm:five}.

\section{Certificates and scope}
\begin{center}
\begin{tabular}{lrrl}
Formula & Variables & Clauses & Producer \\ \hline
Four colours of sizes \(4,4,5,5\) & 242 & 3,438 & Glucose 4 \\
Five colours on \(6\times6\) endpoints & 150 & 87,354 & CaDiCaL 1.9.5 \\
Five colours on \(6\times7\) endpoints & 150 & 146,634 & CaDiCaL 1.9.5
\end{tabular}
\end{center}
The accompanying repository archives the DIMACS formulas and DRAT traces.
Its replay regenerates the clauses before passing them to DRAT-trim.
The three refutations establish the stated finite cases. The unrestricted
conjecture remains open, and we make no priority assertion for the finite
results.

\paragraph{Acknowledgement.}
The search and manuscript used AI assistance. The archived refutations
were replayed with an independent checker.

\begin{thebibliography}{9}
\bibitem{ab} R. Aharoni and E. Berger. Rainbow matchings in \(r\)-partite
\(r\)-graphs. \emph{Electron. J. Combin.} 16 (2009), R119.
\href{https://doi.org/10.37236/208}{doi:10.37236/208}.
\bibitem{dp} M. Delcourt and L. Postle. Finding an almost perfect matching
in a hypergraph avoiding forbidden submatchings. \emph{J. Lond. Math. Soc.}
114 (2026), e70660. \href{https://doi.org/10.1112/jlms.70660}{doi:10.1112/jlms.70660}.
\end{thebibliography}
\end{document}
